\documentclass[11pt]{article}
\usepackage[T1]{fontenc}
\usepackage{lmodern}
\usepackage{amsmath,amssymb,amsthm}
\usepackage{microtype}
\usepackage[margin=1in]{geometry}
\usepackage{booktabs}
\usepackage{enumitem}
\usepackage{xcolor}
\usepackage{listings}
\usepackage{hyperref}

\hypersetup{colorlinks=true,linkcolor=blue!50!black,urlcolor=blue!50!black,citecolor=blue!50!black}

\lstdefinelanguage{lean}{
  keywords={def,theorem,inductive,structure,where,match,with,fun,let,in,namespace,end,open,abbrev,import,by,have,show,exact,termination_by},
  sensitive=true,
  morecomment=[l]{--},
  morecomment=[s]{/-}{-/},
  morestring=[b]",
}
\lstset{language=lean,basicstyle=\ttfamily\small,keywordstyle=\bfseries,
  commentstyle=\itshape\color{black!60},columns=fullflexible,keepspaces=true,
  frame=single,framerule=0.3pt,rulecolor=\color{black!30},xleftmargin=1em,xrightmargin=1em,
  literate={→}{{$\to$}}1 {←}{{$\leftarrow$}}1 {ℕ}{{$\mathbb{N}$}}1 {∀}{{$\forall$}}1
    {≤}{{$\le$}}1 {×}{{$\times$}}1 {⟨}{{$\langle$}}1 {⟩}{{$\rangle$}}1 {λ}{{$\lambda$}}1
    {Λ}{{$\Lambda$}}1 {·}{{$\cdot$}}1 {∃}{{$\exists$}}1 {≫}{{$\gg$}}1}

\newtheorem{theorem}{Theorem}
\newtheorem{lemma}[theorem]{Lemma}
\theoremstyle{definition}
\newtheorem{definition}[theorem]{Definition}
\theoremstyle{remark}
\newtheorem{remark}[theorem]{Remark}

\newcommand{\N}{\mathbb{N}}
\providecommand{\llbracket}{[\![}
\providecommand{\rrbracket}{]\!]}
\newcommand{\lean}[1]{\texttt{#1}}
\newcommand{\BH}{\ensuremath{\mathrm{BH}}}

\title{BusyLean: Naming Ever Larger Numbers in Lean\\[0.3em]
  \large Research notes on constructive large numbers with formal comparisons}
\author{BusyLean contributors}
\date{October 2026}

\begin{document}
\maketitle

\begin{abstract}
We port Cody Roux's \emph{name the biggest number} contest from Coq to Lean~4.
In the contest every number must be a closed, computable definition, and every
claim ``my number is bigger'' must be a kernel-checked theorem. We give an exact
port of the incumbent Contender~5 (the largest numeral of a raw System~T term of
depth~42). We review and rebuild five submitted proposals, and we construct
Contender~6, which combines their best ideas: an object language with W-types
over every finite type, relativized to oracle families, and reflected along
every tree that the language itself can define. Along the way we prove formally
that one proposal's re-implementation of the old language agrees with the exact
port. The notes then explain the strategy for the next contenders, which climb
the ladder System~F, System~F$_\omega$, Calculus of Constructions. Each rung is a
qualitatively stronger total language whose normalization can be justified
inside Lean.
\end{abstract}

\tableofcontents

%======================================================================
\section{The game}

Scott Aaronson's \emph{Who can name the bigger number?} asks for the largest
number one can describe concisely. Cody Roux's repository
\href{https://github.com/codyroux/name-the-biggest-number}{name-the-biggest-number}
adds two caveats that turn the game into one about \emph{proof}. First, numbers
must be \emph{constructive}: computable in principle, with no Busy Beavers.
Second, each new contender must come with a \emph{formal proof} that it beats
the previous one. The original contest uses Coq. Its last entry, Contender~5 by
Samuel Gruetter, is the largest value of a raw STLC+NatRec term (G\"odel's
System~T) of syntax depth at most~42.

\paragraph{The Lean edition.} This repository keeps the contest but switches to
Lean~4 (version 4.34.0). The rules (\texttt{README.md}) are: closed computable
\lean{Nat} definitions; a theorem \lean{contender\_(N-1)\_lt\_\allowbreak contender\_N};
the axioms \lean{propext} and \lean{Quot.sound} allowed and
\lean{Classical.choice} forbidden; at most 15\,s to check the definition and
60\,s to check the proof, from a clean build; and no laziness or shenanigans.
The axiom policy is deliberate. \lean{Classical.choice} is the only one of
Lean's three axioms that conflicts with computability, so it is the one that
must go.

%======================================================================
\section{How to compare unimaginably large numbers}

All numbers here are far too large to evaluate. The formal theorems compare them
directly, but informal \emph{size} is best discussed with the
\emph{fast-growing hierarchy} $f_\alpha$ indexed by ordinals: $f_0(n)=n+1$,
$f_{\alpha+1}(n)=f_\alpha^{n}(n)$, and $f_\lambda(n)=f_{\lambda[n]}(n)$ at
limits. A total language $L$ whose programs are all provably terminating in a
theory $T$ has its bounded-size maxima growing roughly like $f_{|T|}$, where
$|T|$ is the proof-theoretic ordinal of $T$. Two heuristics organize everything
below.

\begin{enumerate}[leftmargin=*]
\item \textbf{The base language dominates.} System~T is calibrated by $\varepsilon_0$.
  T plus one type of countably branching trees, with recursion into all finite
  types, is calibrated by the Bachmann--Howard ordinal
  $\BH=\psi(\varepsilon_{\Omega+1})$. With tree classes at every finite level
  the calibration is $\psi(\Omega_\omega)$, the ordinal of $\Pi^1_1$-CA$_0$.
\item \textbf{Reflection multiplies.} Relativizing a language $L$ to an oracle
  that grows like $f_\beta$ yields growth like $f_{\beta+|L|}$. Iterating this
  along an index of height $\gamma$ yields $f_{|L|\cdot\gamma}$. Explicit
  ordinal notations without collapsing functions stay below \BH, so a stack of
  reflections over a weak base cannot catch up with a stronger base.
\end{enumerate}

These are informal calibrations, not theorems; the formal content of every
contender is a strict inequality checked by Lean's kernel.

%======================================================================
\section{Contender 5, exactly}

\lean{BusyLean/Contender5.lean} ports \lean{Contender.v} one definition at a
time. The port keeps every quirk that can affect the value: raw ill-typed terms
are interpreted with \lean{error} defaults; variables use reversed de~Bruijn
indices; lambdas carry both annotations and cast their body; arrow casts recurse
only when \lean{type\_eqb} succeeds; the enumeration order is the original one,
including right-nested \lean{++}; and the winner is chosen by \lean{maxBy}'s
strict update, starting from \lean{tO}. A smoke test reproduces the reference
table: $0,2,10,146,24976$ terms and maxima $0,0,1,2,3$ at depths $0$--$4$.

\paragraph{Agreement of an independent port.} Proposal P5 re-implemented the old
language \emph{inside} its larger W-language, with W-aware casts and a plain
maximum instead of \lean{maxBy}, and asserted without proof that the result is
Contender~5. We prove it:
\begin{theorem}[\lean{PortAgreement.largest\_eq\_oldMaximum}]
For every depth $n$,
\[\texttt{C5.largestSTLCNatRecNatOfDepth}\ n \;=\; \texttt{BigNumber.oldMaximum}\ n.\]
\end{theorem}
The proof is a logical relation between the two interpreters: equality at
$\N$, preservation at arrows. That avoids function extensionality. Casts in both
directions preserve the relation, including the error branches. The fundamental
lemma covers reversed lookup, binder casts and natural recursion. P5's
enumeration is the translated image of the original, in the same order, and the
\lean{maxBy} winner attains the maximum.

%======================================================================
\section{Five proposals}

Five proposals arrived as archives; \lean{proposals/REVIEW.md} has the details.
All were rebuilt on Lean~4.34.0, and the two that had never been compiled (P4,
P5) needed small repairs, recorded in separate commits.

\begin{center}
\begin{tabular}{llll}
\toprule
& base language & amplifier & informal size\\
\midrule
P1 & System T & FGH over explicit Veblen-style trees & $f_{\varepsilon_0+\alpha}$, $\alpha<\BH$\\
P2 & T + Brouwer trees + oracles & forest reflection on explicit trees & $f_{\BH\cdot\alpha}$\\
P3 & T + unary oracle & reflection on explicit trees & $f_{\varepsilon_0\cdot\alpha}$\\
P4 & T + Brouwer trees + oracle & ramified reflection & $f_{\BH\cdot\alpha}$\\
P5 & T + W-types at all finite types & two diagonals & $f_{\psi(\Omega_\omega)+\omega^2}$\\
\bottomrule
\end{tabular}
\end{center}

P5 has by far the strongest base and the weakest amplifier, which suggests
combining it with the others' amplifiers.

%======================================================================
\section{Contender 6}

\subsection{The object language}
The language is P5's. Types are generated by $A,B::=\N\mid A\to B\mid W\,A$,
where $W\,A$ is the type of well-founded trees branching over~$A$. Terms are raw
(ill-typed terms default), with natural recursion and a W-recursor into every
type. Because $A$ ranges over all types, the language has tree classes of
every finite level, $W\N$, $W(W\N)$, \dots, and its calibration is
$|L|\approx\psi(\Omega_\omega)$.

\subsection{Oracles for free}
P5's interpreter resolves variables by oldest-first levels. Evaluating a program
in the one-element environment $[\langle \N\to\N\to\N, fs\rangle]$ therefore makes
\lean{var 0} denote the oracle family $fs$ everywhere, even under binders. No new
syntax is needed, and the smoke test confirms it: \lean{var0 3 4} evaluates to
$fs(3,4)$, and so does \lean{(fun y => var0 y y) 5} with the family applied at
$(5,5)$.

\subsection{Reflection along the language's own trees}
\begin{lstlisting}
def reflect (fs : Family) (n : ℕ) : ℕ :=
  n + 1 + max (prefixMax (fun i => fs i n) n)
              (maximum (evalWith fs) (allTerms n))
def stage (seed : Family) : WTree ℕ → ℕ → ℕ
  | .leaf   => reflect seed
  | .node f => reflect (fun i => stage seed (f i))
def treesOfDepth (k : ℕ) : WTree ℕ :=
  .node fun i => defTree (nthD (allTerms k) i .zero)
def index : WTree ℕ := .node treesOfDepth
def contender_6 : ℕ := stage p5Family index 42
\end{lstlisting}
At a node the oracle family \emph{is} the family of child stages, so a program
may call any child at a computed index (P2's forest reflection). The new
ingredient is \lean{index}: its $k$-th child joins every $W\N$ tree denoted by a
program of depth $\le k$. The hierarchy therefore runs along the language's own
ordinals, all of $|L|$, instead of an explicit notation system below \BH. By the
multiplication heuristic, Contender~6 is about $f_{|L|\cdot|L|}(42)$.

\subsection{What is proved}
\begin{itemize}[leftmargin=*]
\item \lean{contender\_5\_lt\_contender\_6}: the exact port of Contender~5 is below Contender~6.
\item \lean{p5\_lt\_contender\_6} and \lean{p5Family\_lt\_finalFamily}: P5 is beaten at every input.
\item \lean{defined\_tree\_eventually\_lt}: every tree the language defines is eventually dominated.
\item \lean{program\_lt\_finalFamily}: every program of depth $\le n$, run with the final oracles, is exceeded.
\end{itemize}
All of them depend only on \lean{propext}. The definitions check in about 5\,s
and the proofs in about 6\,s. A machine-checked audit shows that the definition
never mentions \lean{contender\_5} or its winner selection. It also reports,
candidly, that P5's seed contains the old language's maximum as a function of
depth.

%======================================================================
\section{Engineering lessons}

\begin{itemize}[leftmargin=*]
\item \textbf{Toolchain drift changes axiom lists.} In Lean~4.34 a wildcard match
  arm compiles to a \lean{\_sparseCasesOn} helper built on \lean{Nat.land}.
  \lean{Nat.bitwise}'s termination proof uses \lean{propext}, so ``axiom-free on
  4.19'' did not survive the upgrade. The same happens to core \lean{xs[i]?}.
  Exhaustive matches avoid it.
\item \textbf{Tactics carry axioms.} \lean{omega}'s certificates use
  \lean{funext}, and so \lean{Quot.sound}. Two calls were replaced by explicit
  lemma chains.
\item \textbf{Never let the unifier normalize a contender.} Writing
  \lean{iterate f 41 (f 42)} against \lean{iterate f 42 42} made Lean try to
  evaluate the number. Generalizing the statement first, so that only
  $42=41+1$ is checked, fixes it.
\item \textbf{Prove inequalities by inclusion, not monotonicity.} Reflection is
  not monotone in its oracle, because a program may test oracle values. Every
  comparison therefore exhibits an explicit occurrence: a visible oracle index
  or an enumerated program.
\end{itemize}

%======================================================================
\section{Roadmap: impredicativity}

Contender~6 lives below the next collapsing level: everything it does is
predicative. The next qualitative jump is \emph{impredicative polymorphism}.

\paragraph{System F.} The numeric functions definable in Girard's System~F on
Church numerals are exactly those provably total in second-order arithmetic
$\mathsf{Z}_2$. $\psi(\Omega_\omega)$ is the ordinal of $\Pi^1_1$-CA$_0$, a
small fragment of $\mathsf{Z}_2$. A System-F maximum therefore dwarfs
everything above.

\paragraph{Why Lean can host it.} System~F has no set-theoretic model in Lean's
predicative \lean{Type}. Girard's \emph{reducibility candidates} proof of strong
normalization quantifies only over \emph{predicates on terms}, and Lean's
impredicative \lean{Prop} makes $\forall C:\mathrm{Term}\to\mathrm{Prop}$ a
legitimate proposition. The plan:
\begin{enumerate}[leftmargin=*]
\item untyped de~Bruijn $\lambda$-terms with $\beta$-reduction (types erased);
\item Curry-style System~F typing; saturated sets; the interpretation
  $\llbracket\forall X.A\rrbracket_\rho=\bigcap_C\llbracket A\rrbracket_{C\cdot\rho}$;
  and adequacy, which gives strong normalization;
\item Church-style terms (decidable type checking) enumerated by depth, checked
  against the Church numeral type, erased, normalized by a deterministic
  strategy (a total function by well-founded recursion, using the SN proof),
  and decoded;
\item an \emph{oracle constant} $O$ with $\delta$-rule $O\,\overline{n}\to\overline{f(n)}$, so
  System~F programs can amplify a previous contender's family. That makes the
  formal comparison immediate, and it lets F's own strength act on top of it.
\end{enumerate}

\paragraph{Further rungs.} System~F$_\omega$ (higher-order arithmetic) needs
candidates for type constructors, i.e.\ functions between candidate spaces. The
Calculus of Constructions, the system behind Loader's number, adds
dependent types. Each rung is a new contender.

%======================================================================
\section{Status}

\begin{center}
\begin{tabular}{lll}
\toprule
Contender & idea & informal size\\
\midrule
5 (Coq) & System T, depth 42 & $f_{<\varepsilon_0}(42)$\\
6 (Lean) & forest reflection over finite-type W-language, own trees & $f_{\psi(\Omega_\omega)^2}(42)$\\
7 (planned) & System F with oracle & beyond $f_{|\mathsf{Z}_2|}$\\
\bottomrule
\end{tabular}
\end{center}

\end{document}
